\chapter*{General Introduction}
\addcontentsline{toc}{chapter}{General Introduction}

Organizations increasingly depend on connected information systems, cloud services, containers, remote servers, and continuously deployed applications. This evolution increases operational flexibility, but it also increases the number of alerts that security teams must analyze. A Security Operations Center (SOC) receives telemetry from many sources: host detection tools, network intrusion detection engines, runtime security agents, system metrics, application logs, and manual reports. When these signals remain isolated, the analyst must manually connect events, determine impact, decide whether remediation is safe, and verify whether the corrective action actually worked.

The internship project studied in this report addresses this operational problem through \project{}, also named \projectlong{}. The system is a security automation platform that connects alert ingestion, enrichment, incident correlation, AI-assisted investigation, human approval, Ansible remediation, post-fix verification, archiving, monitoring, and dashboard visualization. The project is positioned in the SOAR domain, where orchestration and automation are used to support incident response while preserving analyst control for risky actions.

\section*{Internship Mission}
\addcontentsline{toc}{section}{Internship Mission}

The assigned mission was to design, implement, harden, and validate a backend-centered SOC automation platform able to process alerts from Elasticsearch, normalize them, enrich them with security context, correlate them into incidents, generate response intelligence, and expose the results through an operational web dashboard. The mission also included improvements to runtime remediation, infrastructure monitoring, multi-server asset scoping, API hardening, and validation evidence.

\section*{Problem Statement}
\addcontentsline{toc}{section}{Problem Statement}

The initial problem can be formulated as follows: how can a SOC reduce manual effort and response delay when alerts arrive from heterogeneous sources such as Wazuh, Suricata, Falco, Filebeat, and Telegraf, while still keeping remediation actions controlled, auditable, and verifiable?

Several sub-problems follow from this question. First, raw alert formats differ from one source to another, which makes correlation difficult. Second, repeated or noisy alerts can hide important events. Third, the analyst needs context such as IP information, MITRE ATT\&CK techniques, source history, target host, and related incidents. Fourth, automatic remediation can be dangerous if it executes unvalidated commands or targets the wrong host. Finally, the result of a fix must be verified rather than assumed.

\section*{Objectives}
\addcontentsline{toc}{section}{Objectives}

The project objectives were:

\begin{itemize}[leftmargin=1cm]
\item Implement a multi-source alert pipeline that polls Elasticsearch and normalizes alerts from security and observability sources.
\item Enrich and correlate alerts using GeoIP information, MITRE ATT\&CK mappings, Sigma noise rules, deduplication, and incident grouping.
\item Provide an AI-assisted response layer able to produce investigation summaries, narratives, risk assessments, and Ansible playbooks.
\item Add approval, execution, verification, archiving, audit, and safety mechanisms around remediation workflows.
\item Expose the platform through a FastAPI backend and a Next.js dashboard for SOC analysts and administrators.
\item Validate the implementation through automated tests, compile/build checks, API checks, and runtime verification reports.
\end{itemize}

\section*{Report Organization}
\addcontentsline{toc}{section}{Report Organization}

This report is organized to follow the real engineering story of the project, that is, the order in which the solution was actually built from the ground up. Table~\ref{tab:report-org} maps each chapter to its focus and to the validation evidence that supports it.

\begin{table}[H]
\centering
\caption{Report organization and validation focus}
\label{tab:report-org}
\begin{tabularx}{\textwidth}{@{}p{0.14\textwidth}p{0.48\textwidth}X@{}}
\toprule
\textbf{Chapter} & \textbf{Focus} & \textbf{Validated by} \\
\midrule
Chapter~1 & Host organization, SOC context, existing limits, objectives, methodology & Background review, planning documents \\
Chapter~2 & Huawei Cloud environment and central monitoring stack (ES, Kibana, Wazuh, Suricata, Falco, Filebeat, Telegraf, Redis) & Live ES indices, service health checks \\
Chapter~3 & ARIA design: architecture, data model, pipeline stages, security controls & Architecture diagrams, model and unit tests \\
Chapter~4 & Dashboard functional realization and operational interfaces & Screenshots, API checks, end-to-end scenarios \\
Chapter~5 & Multi-server onboarding, asset scoping, and production deployment & Bootstrap simulation, asset registry, Docker Compose evidence \\
Chapter~6 & Validation, evaluation, limitations, and future work & Compilation, type checks, automated tests, live captures \\
\bottomrule
\end{tabularx}
\end{table}

\textbf{Chapter~1} presents the general context and project scope. \textbf{Chapter~2} builds the telemetry foundation. \textbf{Chapter~3} explains the design and internal logic of ARIA. \textbf{Chapter~4} shows the functional realization through the dashboard. \textbf{Chapter~5} covers onboarding and production deployment. \textbf{Chapter~6} validates the solution and discusses limitations and future work. The report ends with a general conclusion and appendices containing technical inventories, configuration extracts, and validation evidence.

